\documentclass{leveldoc}
\DocUses{foundations,first_bend,slingshot,around_the_wall,einzel_lens}
\DocFingerprint{1fa34953959d976d}

\begin{document}

\section{What the level asks}\label{sec:asks}

The game's text:
\leveldescription

Two ions ($q = \val{q}$, $m = \val{m}$) leave $(\val{xLaunch}, \val{yLaunch})$ along $+x$,
with kinetic energies $T = \val{Tslow}$ and $T = \val{Tfast}$. On both shots the detector
is the rectangle $\val{xDlo} \le x \le \val{xDhi}$, $\val{yDlo} \le y \le \val{yDhi}$,
beside the source and below it. An entry counts at every direction and every kinetic
energy. The player places at most $\val{maxCharges}$ charges, of the particle's sign
only, inside the region $\val{xRegLo} \le x \le \val{xRegHi}$,
$\val{yRegLo} \le y \le \val{yRegHi}$, and no magnets. A flight ends when the particle
touches a charge, leaves the world, enters the detector or runs out of time
(\xsecref{foundations}{sec:board}); $t_{\max} = \val{tmax}$.

The reference solution adds $Q = \val{Qlo}$ at $(\val{xLo}, \val{yLo})$ and
$Q = \val{Qhi}$ at $(\val{xHi}, \val{yHi})$. Both have the particle's sign, so they
stand in front of the ions and repel. Each ion slows, turns, and comes back into the
detector: the slower one through the right face, the faster one through the top face.
The level asks for those two arrivals. This document integrates the two-charge flights,
and the flight of each charge alone, by the fixed-step method of
\levelref{around_the_wall}.

\section{Setup}\label{sec:setup}

The level file, field by field (the appendix reproduces it):
\setuptable

In symbols: $c = \val{c}$, the particle has radius $a = \val{aRad}$, and every charge
has radius $R = \val{R}$. The two launches are
\[
  \gamma = 1 + \frac{T}{mc^2}, \qquad
  \abs{\vect{p}_0} = \sqrt{\frac{T^2}{c^2} + 2mT}, \qquad
  v_0 = \frac{\abs{\vect{p}_0}}{\gamma m}.
\]
For $T = \val{Tslow}$ this is $\gamma = \val{gammaSlow}$,
$\abs{\vect{p}_0} = \val{pSlow}$ and $v_0 = \val{vSlow}$. For $T = \val{Tfast}$ it is
$\gamma = \val{gammaFast}$, $\abs{\vect{p}_0} = \val{pFast}$ and $v_0 = \val{vFast}$.
The direction in the file is $+x$ on both shots. The world is $-1 \le x \le \val{xWorld}$,
$-1 \le y \le \val{yWorld}$. The couplings are $\kappa_\mathrm{l} = \val{kLo}$ and
$\kappa_\mathrm{h} = \val{kHi}$, in the file's order.

\begin{figure}[htbp]
\centering
\begin{tikzpicture}
\begin{axis}[axis equal image, width=0.88\linewidth, xmin=-1.8, xmax=32.4, ymin=-2.2, ymax=22.2,
  xtick={0,5,...,30}, ytick={0,5,...,20}, tick label style={font=\small},
  xlabel={$x$}, ylabel={$y$}, label style={font=\small}]
  \draw[black!45, dashed] (-1,-1) rectangle (31,21);
  \draw[black!50] (0,0) rectangle (30,20);
  \draw[black!40, dashed] (18,0) rectangle (29,20);
  \fill[green!50!black, fill opacity=0.18] (0,4) rectangle (2,7);
  \draw[green!40!black] (0,4) rectangle (2,7);
  \fill[red!70!black, fill opacity=0.35] (\rawval{xLo},\rawval{yLo}) circle[radius=0.3];
  \draw[red!70!black] (\rawval{xLo},\rawval{yLo}) circle[radius=0.3];
  \fill[red!70!black, fill opacity=0.35] (\rawval{xHi},\rawval{yHi}) circle[radius=0.3];
  \draw[red!70!black] (\rawval{xHi},\rawval{yHi}) circle[radius=0.3];
  \addplot[thin, black!55] table {pathlo.dat};
  \addplot[thin, purple!70!black] table {pathhi.dat};
  \addplot[very thick, blue!70!black] table {pathslow.dat};
  \addplot[thick, orange!80!black] table {pathfast.dat};
  \draw[-{Stealth[length=6pt]}, thick] (\rawval{xLaunch},\rawval{yLaunch}) --
    (\rawval{arrowX},\rawval{arrowY});
  \fill (\rawval{xLaunch},\rawval{yLaunch}) circle[radius=0.11];
  \node[font=\scriptsize, red!60!black, anchor=north] at (18,10.55) {$Q_\mathrm{l}$};
  \node[font=\scriptsize, red!60!black, anchor=south] at (21,17.45) {$Q_\mathrm{h}$};
  \node[font=\scriptsize, green!30!black, anchor=east] at (-0.15,5.5) {detector};
\end{axis}
\end{tikzpicture}
\caption{The reference ions turn and enter the detector. Blue is the slower ion, orange
the faster one. Grey is the lower charge alone, on the slower ion: it leaves through the
top of the world. Purple is the upper charge alone, on the slower ion: it leaves through
the floor. The discs are the two repelling charges, each of radius $R$. The dashed
rectangle is the region in which the player may place charges.}
\label{fig:board}
\end{figure}

\section{Principles and model}\label{sec:model}

\begin{principles}
  \item The field of a charged sphere and the force on a non-polarizable particle,
    \xeqref{foundations}{eq:coulomb} and \xeqref{foundations}{eq:sphereforce}.
  \item The relativistic equation of motion \xeqref{foundations}{eq:eom}, the launch
    \xeqref{foundations}{eq:launch} and the energy \xeqref{foundations}{eq:energy}. The
    motion stays in the plane (\xsecref{foundations}{sec:plane}).
  \item The board's rules (\xsecref{foundations}{sec:board}). The first entry into the
    detector decides. This level adds no direction and no kinetic window.
  \item The classical Runge--Kutta method \xeqref{foundations}{eq:rk4} and the fixed-step
    flight of \xsecref{around_the_wall}{sec:method}. With two centres the orbit equation
    does not separate (\xeqref{slingshot}{eq:noseparation}).
  \item The radiated power \xeqref{foundations}{eq:larmor}, with the acceleration
    \xeqref{first_bend}{eq:acceleration}, summed by Simpson's rule
    \xeqref{foundations}{eq:simpson}.
\end{principles}

The level has no magnets, so $\vect{B} = \vect{0}$. The potential and the kinetic energy
are \xeqref{einzel_lens}{eq:potential} and \xeqref{einzel_lens}{eq:kinetic}, here with
$\kappa_\mathrm{l} = \val{kLo}$ at $(\val{xLo}, \val{yLo})$ and
$\kappa_\mathrm{h} = \val{kHi}$ at $(\val{xHi}, \val{yHi})$. The force is
\xeqref{around_the_wall}{eq:force} with $N = \val{nCharges}$, the charges in the file's
order. The state $(x, y, p_x, p_y)$ obeys \xeqref{foundations}{eq:eom} with that force
until the first event.

\begin{assumption}{The charges are held in place}\label{as:fixed}
None of the charges recoils. That is a rule of the game
(\xsecref{foundations}{sec:radiation}).
\end{assumption}

\begin{assumption}{No radiation}\label{as:radiation}
The radiated energy is not taken from the particle. Over the six flights of this
document the largest fraction of the launch energy is $\val{radFracMax}$
(§\ref{sec:hand}). The game makes the same approximation and reports the same integral
(\physref{8}).
\end{assumption}

\begin{assumption}{A point particle}\label{as:point}
$a = 0$, and the particle does not polarize, so \xeqref{foundations}{eq:sphereforce} is
exact. On the two reference flights the closest approach to either charge is
$\val{rClose}$, on the \val{rCloseShot} ion, to the \val{rCloseCharge} charge. The
clearance outside the sphere is $\val{clearClose}$. The samples that locate it are at
most $\val{sampleArc}$ apart (§\ref{sec:method}), far below that clearance.
\end{assumption}

\begin{assumption}{The step resolves the turn}\label{as:step}
The reported step is $h = \val{H}$. Halving it estimates the largest error of an end
state as $\val{errMax}$ (§\ref{sec:hand}). That is far below the height of the
detector, from $\val{yDlo}$ to $\val{yDhi}$. The tolerance used in
§\ref{sec:comparison} is the larger of $10^{-9}$ and ten times this estimate.
\end{assumption}

\section{Derivation}\label{sec:derivation}

\subsection{What the energy allows}

At the launch the potential is $V(\vect{x}_0) = \val{VZero}$ for both ions, since they
leave the same point. Then
\[
  W - mc^2 = T + V(\vect{x}_0)
\]
is $\val{WbindSlow}$ for the slower ion and $\val{WbindFast}$ for the faster one. Both
are positive. Far from the charges $V$ falls to $0$, which lies below either bound, so
\xeqref{einzel_lens}{eq:kinetic} allows an ion to leave the world. The force at the
launch is $(\val{FxZero}, \val{FyZero})$. The horizontal part is negative: both charges
stand ahead of the ion and repel, so the mirror retards the beam. The lower charge sits
below the launch and pushes upward; the upper charge sits above it and pushes downward.
The upper coupling is the larger one, and the sum $\val{FyZero}$ points downward.

Both couplings are positive, so $V$ rises as an ion approaches the charges and
\xeqref{einzel_lens}{eq:kinetic} falls. On the reference flights the kinetic energy
reaches $\val{TminSlow}$ and $\val{TminFast}$, far below the launch values, and the
horizontal momentum changes sign. That is the mirror. At the detector the potential has
fallen below its value at the launch, and the arrival kinetic energies
$\val{TarrSlow}$ and $\val{TarrFast}$ lie above $\val{Tslow}$ and $\val{Tfast}$.

\subsection{Why the flight is integrated}\label{sec:noseparation}

Two fixed centres fail to separate in the relativistic problem
(\xeqref{slingshot}{eq:noseparation}), and this pair is staggered, so the obstruction
stands. Each ion is the initial-value problem for \xeqref{around_the_wall}{eq:force}
at $N = \val{nCharges}$. Section~\ref{sec:method} computes it to a stated accuracy.

\subsection{The method}\label{sec:method}

The state is advanced by the classical Runge--Kutta method \xeqref{foundations}{eq:rk4}
at the fixed step $h$, and a boundary is crossed on the cubic Hermite interpolant
\xeqref{slingshot}{eq:hermite}, as specified in \xsecref{around_the_wall}{sec:method}.
After each step the interpolant is sampled at $\val{nSample}$ points. A sphere, the
world's edge or the detector is tested only when the box of the step's ends, enlarged
by $ch$, meets it. The samples are read in the board's order
(\xsecref{foundations}{sec:board}). The first instant of the winning event is the
bisection of the interpolant, $\val{nBisect}$ iterations. The position reported is the
interpolant at that instant, and the face coordinate is then set to the face value.
The momentum is one further step of the method, of that length. The error of the
reported step is a fifteenth of its difference from the step $2h$
(\xsecref{foundations}{sec:numerics}).

On each step the same interpolant is also sampled at nine equally spaced fractions.
\xeqref{einzel_lens}{eq:kinetic} at those samples gives the kinetic-energy dip, the
distances to the two centres give the closest approach, and the greatest abscissa gives
the turn. The samples are an eighth of a step apart in time, so a point between them
lies within $\val{sampleArc}$ of a sample. Those three quantities are located to that
tolerance. The end state is held to the step-halving estimate.

\section{The reference solution by hand}\label{sec:hand}

\begin{procedure}
  \item\label{st:launch} \emph{The launches} (\xeqref{foundations}{eq:launch}):
    the slower ion has $\gamma = \val{gammaSlow}$ and $\abs{\vect{p}_0} = \val{pSlow}$;
    the faster ion has $\gamma = \val{gammaFast}$ and $\abs{\vect{p}_0} = \val{pFast}$.
    Both are aimed along $+x$. The force is $(\val{FxZero}, \val{FyZero})$, retarding
    and downward.
  \item\label{st:energy} \emph{The energy}: $W = \val{WSlow}$ and
    $W - mc^2 = \val{WbindSlow}$ for the slower ion, and $W = \val{WFast}$ and
    $W - mc^2 = \val{WbindFast}$ for the faster one. For both ions the world's
    edge is energetically open.
  \item\label{st:flight} \emph{The slower ion}: the method of §\ref{sec:method} with
    $h = \val{H}$. The states every five time units, and the drift of the energy:
    \begin{center}\small\val{states}\end{center}
    The kinetic energy falls to $\val{TminSlow}$ at $t = \val{tTmin}$, at
    $(\val{xTmin}, \val{yTmin})$. The rightmost sample is $(\val{xTurnSlow},
    \val{yTurnSlow})$ at $t = \val{tTurnSlow}$. The ion has almost stopped between the
    two charges, and the horizontal momentum in the table changes sign there.
  \item\label{st:entry} \emph{The two entries.} The slower ion's step from
    $t_n = \val{tBefore}$, at $(\val{xBefore}, \val{yBefore})$, ends at
    $(\val{xAfter}, \val{yAfter})$. The interpolant crosses the \val{faceE} face at the
    fraction $u = \val{uE}$. The residual before the face coordinate is assigned is
    $\val{faceRes}$, so $x_\mathrm{e} = \val{xE}$, $y_\mathrm{e} = \val{yE}$ and
    $t_\mathrm{e} = \val{tE}$, after $\val{steps}$ steps. One step of length
    $t_\mathrm{e} - t_n$ gives $\vect{p}_\mathrm{e} = (\val{pxE}, \val{pyE})$, at
    $\val{angE}^\circ$ to $+x$. The horizontal part is negative: the ion is coming back.

    The faster ion's kinetic energy falls to $\val{TminFast}$, and its rightmost sample
    is $(\val{xTurnFast}, \val{yTurnFast})$ at $t = \val{tTurnFast}$, farther in than
    the slower ion's turn. It then enters the \val{RefFastFace} face at
    $(\val{RefFastX}, \val{RefFastY})$, at $t = \val{RefFastT}$, with momentum
    $(\val{RefFastPx}, \val{RefFastPy})$, at $\val{RefFastAng}^\circ$.

    The slower ordinate lies between $\val{yDlo}$ and $\val{yDhi}$. The faster abscissa
    lies between $\val{xDlo}$ and $\val{xDhi}$.
  \item\label{st:check} \emph{The checks}: along the slower flight the energy stays
    within $\val{wDev}$ of $W$. The arrival kinetic energies from
    \xeqref{einzel_lens}{eq:kinetic} are $\val{TarrSlow}$ and $\val{TarrFast}$. Halving
    the step estimates the two end states as $\val{RefSlowErr}$ and $\val{RefFastErr}$.
    The same outcomes appear at $h$ and at $2h$.
  \item\label{st:rad} \emph{The radiation}: the power \xeqref{foundations}{eq:larmor}
    with the acceleration \xeqref{first_bend}{eq:acceleration}, summed by Simpson's rule
    on the steps of $h/2$. The slower ion radiates $\val{RefSlowRad}$ and the faster
    ion $\val{RefFastRad}$. The relative differences from the same sums at step $h$,
    divided by fifteen, are $\val{RefSlowRadErr}$ and $\val{RefFastRadErr}$.
\end{procedure}

\begin{result}[Both ions return]
The slower ion enters the detector through its right face at $(\val{xE}, \val{yE})$ at
$t = \val{tE}$, coming back. The faster ion enters through the top face at
$(\val{RefFastX}, \val{RefFastY})$ at $t = \val{RefFastT}$, also coming back. The level
is solved.
\end{result}

By hand, Step~\ref{st:flight} is $\val{steps}$ steps of four evaluations of the force,
for the slower ion, and the same method for the faster one. Across all six flights the
largest step-halving estimate is $\val{errMax}$.

\section{Comparison with the game}\label{sec:comparison}

\comparisontable

The six outcomes agree. The tolerance on a state is $\val{tolFlight}$. Ten times the
largest step-halving estimate $\val{errMax}$ lies below $10^{-9}$, so the tolerance is
that floor. The radiated energies are compared relatively, with tolerance $10^{-4}$.
The state tolerance and the radiation tolerance were fixed from these estimates before
the comparison.

The largest absolute difference is the slower ion's arrival time,
$\val{cmpRefSlowTDiff}$. Its ordinate differs by $\val{cmpRefSlowYDiff}$. The
difference is many times $\val{errMax}$ and a small fraction of the tolerance. The
game's verified end states are good to about $10^{-10}$
(\xsecref{foundations}{sec:verification}), so a difference of this size is what that
accuracy allows when two centres are summed.

The faster ion's ordinate agrees exactly, and so do both ordinates of the lower charge
alone. The slower ion's abscissa differs by $\val{cmpRefSlowXDiff}$.

The quadrature uncertainties of the two reference ions are $\val{RefSlowRadErr}$ and
$\val{RefFastRadErr}$, inside the relative tolerance. The largest relative difference
is the lower charge's faster ion, $\val{cmpLoFastRadDiff}$.

\section{Other solutions}\label{sec:others}

Either added charge, without the other, sends both ions out of the world. The method
of §\ref{sec:method} integrates both (\code{placements.json}), each ion at
$h = \val{H}$ and at $2h$.
\begin{enumerate}[label=\arabic*.]
  \item $Q = \val{Qlo}$ at $(\val{xLo}, \val{yLo})$ only. The charge sits below the
    launch, so it pushes the beam upward. Both ions leave through the top face of the
    world.

    The slower ion leaves at $(\val{LoSlowX}, \val{LoSlowY})$, at $t = \val{LoSlowT}$.

    The momentum is $(\val{LoSlowPx}, \val{LoSlowPy})$, at $\val{LoSlowAng}^\circ$ to
    $+x$. The horizontal part is still positive. The estimate is $\val{LoSlowErr}$.
    The radiated energy is $\val{LoSlowRad}$.

    The faster ion leaves at $(\val{LoFastX}, \val{LoFastY})$, at $t = \val{LoFastT}$.

    The momentum is $(\val{LoFastPx}, \val{LoFastPy})$, at $\val{LoFastAng}^\circ$.
    The estimate is $\val{LoFastErr}$. The radiated energy is $\val{LoFastRad}$.
  \item $Q = \val{Qhi}$ at $(\val{xHi}, \val{yHi})$ only. The charge sits above the
    launch, so it pushes the beam downward. Both ions leave through the bottom face of
    the world.

    The slower ion leaves at $(\val{HiSlowX}, \val{HiSlowY})$, at $t = \val{HiSlowT}$.

    The momentum is $(\val{HiSlowPx}, \val{HiSlowPy})$, at $\val{HiSlowAng}^\circ$.
    The estimate is $\val{HiSlowErr}$. The radiated energy is $\val{HiSlowRad}$.

    The faster ion leaves at $(\val{HiFastX}, \val{HiFastY})$, at $t = \val{HiFastT}$.

    The momentum is $(\val{HiFastPx}, \val{HiFastPy})$, at $\val{HiFastAng}^\circ$.
    The estimate is $\val{HiFastErr}$. The radiated energy is $\val{HiFastRad}$.
\end{enumerate}

The reference needs both charges. The lower one supplies the upward push that keeps the
beam off the floor; the upper one supplies the downward push that keeps it off the
ceiling. Together the beam slows almost to a stop (Step~\ref{st:flight}) and returns
to the detector, as in Figure~\ref{fig:board}. Section~\ref{sec:comparison} records the
outcomes.

\section{What the game computes}\label{sec:engine}

The game integrates \xeqref{foundations}{eq:eom} with its eighth-order Runge--Kutta
method at the tolerances $10^{-10}$ and $10^{-12}$ (\physref{5}). It finds the crossing
on its interpolant (\physref{6}) and takes the first entry into the box as the outcome
(\xsecref{foundations}{sec:board}). It verifies that outcome by comparing its two runs
(\physref{7}). The model is this document's, including \ref{as:radiation} (\physref{8}).
The agreement in §\ref{sec:comparison} is between the fixed-step flight of
§\ref{sec:method} and that integration.

\section{Further reading}\label{sec:reading}

\begin{itemize}
  \item \levelref{einzel_lens}: the two-charge potential and the kinetic energy, used
    here for a mirror instead of a lens that the rays pass through.
  \item \levelref{around_the_wall}: the force of several centres and the fixed-step
    method with the enlarged-box test.
  \item \levelref{slingshot}: why two fixed centres have no closed relativistic orbit.
  \item \levelref{first_bend}: the acceleration that enters the radiated power.
  \item \xsecref{foundations}{sec:board} and \xsecref{foundations}{sec:numerics}: the
    order of events, and the step-halving estimate.
\end{itemize}

\levelfileappendix

\end{document}
